% 中文论文骨架。正式编译：uv run --locked python -m scripts.paper_template main.tex --contest cumcm --compile --output-directory outputs/build-001
% 全国赛论文格式规范（2026年修订稿，2026-10-09核对）：A4、四边页边距不少于 2.5 厘米、页脚居中的阿拉伯页码（从摘要页起）、
% 摘要页含标题和关键词且不超过一页、正文不要目录、不超过30页（摘要单页、附录不计入正文）、附录页数不限且含支撑材料文件清单和完整可运行源程序
% （没有程序要写明“本论文没有用到程序”）、论文中不得出现身份、学校和赛区信息。字体与字号不做统一要求。
% 编译后检查：uv run --locked python -m scripts.contest_rules main.pdf --contest cumcm --forbidden 学校名。以当年官方文件为准。
\newcommand{\PraxisTitle}{论文题目}
% BEGIN PRAXIS CUMCM STYLE v2
% Praxis CUMCM house style v2; content and contest rules are separate.
\documentclass[a4paper,12pt,fontset=none]{ctexart}
\usepackage[top=2.54cm,bottom=2.54cm,left=2.8cm,right=2.8cm]{geometry}
\setCJKmainfont[BoldFont=FandolSong-Bold.otf,ItalicFont=FandolSong-Regular.otf,BoldItalicFont=FandolSong-Bold.otf]{FandolSong-Regular.otf}
\setCJKsansfont[BoldFont=FandolHei-Bold.otf]{FandolHei-Regular.otf}
\setCJKmonofont[BoldFont=FandolSong-Bold.otf]{FandolSong-Regular.otf}
\usepackage{amsmath,amssymb}
\usepackage{newtxtext,newtxmath}
\usepackage{booktabs,graphicx,xcolor,pgfplots,caption,needspace,fvextra,array,tabularx}
\usepackage[hidelinks,pdfauthor={},pdftitle={\PraxisTitle}]{hyperref}
\urlstyle{same}\Urlmuskip=0mu plus 1mu
\renewcommand{\CJKttdefault}{\CJKrmdefault}
\pgfplotsset{compat=1.17}
\usepgfplotslibrary{groupplots}
\definecolor{muted}{HTML}{7A7F85}
\definecolor{main}{HTML}{0072B2}\definecolor{accent}{HTML}{D55E00}
\linespread{1.3}
\setlength{\parskip}{0pt}
\ctexset{section={format=\bfseries\zihao{4}\centering,name={,},number=\arabic{section},aftername=\quad,beforeskip=2.4ex plus .4ex,afterskip=1.6ex},
 subsection={format=\bfseries\zihao{-4}\raggedright,name={,},number=\arabic{section}.\arabic{subsection},aftername=\quad,beforeskip=1.8ex plus .3ex,afterskip=1ex}}
\DeclareCaptionLabelFormat{heiti}{{\bfseries #1\,#2}}
\captionsetup{labelformat=heiti,labelsep=quad,font=small,justification=centering}
\captionsetup[table]{position=above,skip=5pt}\captionsetup[figure]{position=below,skip=4pt}
\setlength{\abovedisplayskip}{8pt}\setlength{\belowdisplayskip}{8pt}
\pagestyle{plain}
% END PRAXIS CUMCM STYLE
\begin{document}
\begin{center}{\bfseries\zihao{3} 论文题目：写出结论，而不是只写主题}\end{center}
\section*{摘\hspace{1em}要}
最后写本页：按问题一、问题二……逐问给出方法、关键数值（含单位）、检验与结论，不超过一页。\par
关键词：关键词一；关键词二

\clearpage
\section{问题重述与数据核对}
\section{假设、定义与符号}
\begin{table}[htbp]\centering\small\caption{符号与单位}
\begin{tabular}{ll}\toprule 符号 & 含义与单位\\ \midrule $x_i$ & 示例变量，元\\ \bottomrule\end{tabular}\end{table}
\section{模型与求解}
\begin{equation}R(\mathbf{x})=\max_{i}\frac{q_ix_i}{M}\le\rho\end{equation}
\begin{figure}[htbp]\centering
\begin{tikzpicture}
\begin{axis}[width=0.9\linewidth,height=6cm,xlabel={风险上限（\%）},ylabel={净收益率（\%）},grid=major]
\addplot[color=main,line width=1.1pt,no marks] coordinates {(0,5) (1,20) (2,25)};
\end{axis}
\end{tikzpicture}
\caption{图题用一句话说明读者应看到什么}\end{figure}
\section{检验、敏感性与局限}
与简单基线比较；扰动参数和执行误差；写明模型不能支持的结论。
\section{结论与建议}

\section*{参考文献}
\begingroup\small\sloppy\raggedright\setlength{\parindent}{0pt}\hangindent=2em
[1] 作者. 题名[J]. 刊名, 年, 卷(期): 页码.\par
\endgroup
\clearpage\section*{附录 A 支撑材料文件列表}
% 列出支撑材料压缩包内的全部文件；确实没有支撑材料时写“本论文没有支撑材料”。
\clearpage\section*{附录 B 完整源程序}
\begin{Verbatim}[fontsize=\scriptsize,breaklines,breakanywhere,numbers=left,numbersep=5pt,xleftmargin=1.6em]
# 粘贴完整可运行源程序；确实没有用到程序时删去本节并写“本论文没有用到程序”
\end{Verbatim}
\end{document}
